\subsection{Modelowanie relacji przestrzennych w warunkach niepełnych danych}

\paragraph{}Istotnym obszarem badań w dziedzinie rozpoznawania twarzy jest opracowanie metod odpornych na degradację lub zmianę cech biometrycznych, do czego dochodzi zarówno w przypadku okluzji (przysłonięcia), jak i naturalnych procesów starzenia. Współczesna literatura wskazuje, że klasyczne sieci splotowe (CNN) są skuteczne w ekstrakcji cech lokalnych, jednak mogą napotykać trudności w modelowaniu globalnych zależności geometrycznych, zwłaszcza gdy część danych wejściowych ulega degradacji lub zasłonięciu.

W tym kontekście obiecującym kierunkiem jest wykorzystanie sieci typu Transformer do modelowania długozasięgowych zależności przestrzennych (ang. long-range dependencies). Jak wykazują Prakash i Rai [2], relacje przestrzenne pomiędzy kluczowymi elementami twarzy pozostają stosunkowo stałe, nawet gdy powierzchowne cechy (tekstura skóry, zmarszczki, proporcje) ulegają znacznym zmianom w czasie. Autorzy ci proponują hybrydowe podejście, w którym Transformer nie zastępuje głównej sieci ekstrakcyjnej (backbone CNN), lecz wykorzystywany jest jako pomocnicza funkcja straty (ang. auxiliary loss), współtworząc mechanizm Transformer-Enhanced Loss.

\paragraph{}W proponowanej przez nich architekturze, mapa cech z ostatniej warstwy splotowej jest traktowana jako sekwencja wektorów kontekstowych. Mechanizm atencji pozwala sieci skupić się na zachowanych relacjach między widocznymi fragmentami twarzy, co mityguje wpływ zakłóceń. W badaniach na zbiorach o dużej wariancji wiekowej (CA-LFW, AgeDB), połączenie standardowej straty metrycznej (np. ArcFace) ze stratą transformatorową pozwoliło na osiągnięcie wyników typu State-of-the-Art [2].

Dla problematyki okluzji wnioski te są kluczowe: sugerują, że skuteczna identyfikacja twarzy częściowo przysłoniętej nie musi polegać wyłącznie na rekonstrukcji brakujących fragmentów (inpainting), lecz może opierać się na silniejszym modelowaniu wzajemnych relacji przestrzennych tych fragmentów, które pozostają widoczne. Wykorzystanie sekwencyjnej analizy map cech pozwala bowiem na utrzymanie spójności reprezentacji tożsamości pomimo braku ciągłości wizualnej w obrazie wejściowym.